\section{Terminologi Dasar}
Untuk memahami kriptografi, kita perlu mengenal beberapa istilah kunci:
\begin{itemize}
    \item \textbf{Plainteks (\textit{plaintext})}: Pesan asli yang dapat dibaca dan dimengerti maknanya.
    \item \textbf{Cipherteks (\textit{ciphertext})}: Pesan yang telah disandikan sehingga tidak bermakna bagi pihak yang tidak berhak.
    \item \textbf{Enkripsi (\textit{encryption})}: Proses mengubah plainteks menjadi cipherteks.
    \item \textbf{Dekripsi (\textit{decryption})}: Proses mengembalikan cipherteks menjadi plainteks.
    \item \textbf{Kunci (\textit{key})}: Parameter rahasia yang digunakan dalam proses enkripsi dan dekripsi.
    \item \textbf{Cipher}: Algoritma atau aturan yang digunakan untuk enkripsi dan dekripsi.
\end{itemize}

\subsection{Sistem Kriptografi}
Secara formal, sebuah sistem kriptografi dapat dinyatakan sebagai \textit{quintuple} $(\mathcal{E},\allowbreak \mathcal{D},\allowbreak \mathcal{M},\allowbreak \mathcal{K},\allowbreak \mathcal{C})$ di mana:
\begin{itemize}
    \item $\mathcal{M}$ adalah himpunan plainteks.
    \item $\mathcal{K}$ adalah himpunan kunci.
    \item $\mathcal{C}$ adalah himpunan cipherteks.
    \item $\mathcal{E}$ adalah himpunan fungsi enkripsi $E: \mathcal{M} \times \mathcal{K} \to \mathcal{C}$.
    \item $\mathcal{D}$ adalah himpunan fungsi dekripsi $D: \mathcal{C} \times \mathcal{K} \to \mathcal{M}$.
\end{itemize}

